\documentclass[../00_Main.tex]{subfiles}
\begin{document}

Funding constraints regarding the manufacturing of the conductive cranial shell led to an adaptation of the validation strategy, which originally aimed for a full physical integration of the software with the anatomical phantom head. Consequently, the system evaluation was conducted using a modular verification framework. Instead of a single final assembly test, each subsystem was isolated and tested independently to ensure strict compliance with the functional requirements:
\begin{itemize}
    \item The \textbf{Software Subsystem} was evaluated for latency and waveform spectral purity.
    \item The \textbf{Hardware Transduction} was validated using gelatine cubes and controlled input signals.
    \item The \textbf{Signal Acquisition Subsystem} was tested for signal integrity and noise levels using controlled input signals.
    \item The \textbf{Closed-Loop System} was assessed by integrating the signal generation and a simplified form of the hardware module, comparing the output using the acquisition module to verify real-time performance.
\end{itemize}

The validation process utilized both laboratory instrumentation and the developed software subsystem to generate controlled input signals.

To validate the behavior of the complete integrated system in the absence of the final physical assembly with the anatomical phantom head, the project utilizes
the \textbf{Phantom Head Numerical Model} as the integration reference for a two-layer phantom head constructed following the methods described in 
Section~\ref{sec:materials_and_methods}. This theoretical validation demonstrates that, given the proven performance of the individual components, the complete system functions according to the design specifications when the modules are combined.

\subsection{Materials: Electrical Behavior and Suitability for Phantom Construction}

The phantom head design relies on a bilayer architecture composed of an external conductive
shell and an internal conductive volume. These two layers are subject to fundamentally
different material constraints. The outer shell is designed to be fabricated using a commercially available conductive
polymer, selected to balance electrical performance, mechanical printability, and practical
accessibility within an academic prototyping context. While custom composite formulations
or fully bespoke materials could, in principle, be developed to achieve tighter control
over electrical properties, such approaches would significantly increase fabrication
complexity and cost. For this reason, a commercially available conductive PLA was adopted
as a representative and reproducible choice, with electrical properties largely determined
by the material specification and the additive manufacturing process itself. Accordingly,
its characterization relies primarily on manufacturer data and existing literature rather
than on experimental tunability.

In contrast, the inner conductive layer is based on a saline--gelatine medium, which is
readily available, inexpensive, and highly customizable through formulation parameters.
For this reason, the experimental characterization efforts focused on the inner layer, with
the goal of assessing its electrical behavior, tunability, and reproducibility in the
context of phantom construction. Rather than aiming at a precise metrological determination
of conductivity values, the experiments were designed to qualitatively evaluate how the
electrical response of the material evolves with ionic concentration, and whether these
properties can be reliably reproduced across fabrication batches.

\subsubsection{Gelatine Layer}
\paragraph*{Impedance Analysis and Material Sensitivity}

The data reveals a pronounced inverse relationship between NaCl concentration and the magnitude of electrical impedance ($|Z|$) across the tested frequency spectrum. As the NaCl concentration increases, the impedance decreases, indicating an increase in conductivity. This trend is consistent with the expected behavior of ionic solutions, where higher salt content facilitates charge transport. We also observe a decrease in the magnitude of the impedance as the frequency increases. 

However, the addition of NaCl and increase of conductivity introduces a complex trade-off regarding the phase response. While increasing the NaCl concentration to 5\% effectively lowers the resistive component, the data indicates that it simultaneously induces a more capacitive response compared to the 0\% (or low-concentration) baseline.

Specifically, the phase angle ($\theta$) for the solution without NaCl remains closer to 0 degrees, indicating a predominantly resistive behavior. In contrast, the 5\% formulation, while lower than the 2.5\% solution, exhibits a larger phase shift. This increased capacitive reactance is attributable to interfacial polarization. The higher density of free ions facilitates the formation of a stronger electrical double-layer (EDL) at the electrode-gelatine interface, thereby increasing the effective capacitance of the system alongside its conductivity, especially at lower frequencies.

These findings align with the results obtained by Owda \emph{et al.}~\cite{Owda2020} in their study of gelatine-based phantoms, where they also reported similar trends in impedance and phase behavior with varying salt concentrations.

The reproducibility analysis between two batches of the 2.5\% NaCl gelatine samples demonstrated a high degree of consistency in both impedance magnitude and phase response. The negligible variance observed between batches indicates that the manufacturing process is reliable and produces uniform material properties, which is crucial for phantom applications requiring repeatable electrical characteristics.

\paragraph*{Signal Localization}

The spatial distribution of the measured voltages confirms the effective propagation and localization of the dipole signal within the conductive medium, as getting closer to each source results in a higher measured potential at the corresponding electrode frequency.

\subsection{Numerical Phantom Validation and System-Level Implications}

The numerical modeling strategy adopted in this work was validated in a staged manner,
progressing from analytically tractable benchmark problems to the full phantom head
geometry. Initial validation was performed using simplified spherical domains for which
closed-form or well-established reference solutions exist. These benchmarks primarily
quantify discretization error and numerical consistency, allowing the finite element
implementation to be verified independently of geometric complexity. Across all tested
configurations, reconstruction errors were consistently low and the numerical behavior
remained stable, providing confidence in the correctness of the governing equations,
boundary condition implementation, and solution strategy.

While these spherical benchmarks establish numerical accuracy, they do not capture the
geometric complexity of a realistic head-shaped phantom. Unlike the sphere, the phantom
head geometry exhibits nonuniform curvature, localized geometric features, and a lack of
global symmetry, all of which can influence field distributions and numerical sensitivity.
Additional modeling discrepancies may also arise from idealizations in the material
representation, including the assumption of homogeneous layers and the omission of
manufacturing-induced microstructures.

In particular, more detailed physical modeling could be achieved by further mesh
refinement or by explicitly representing features introduced by the additive manufacturing
process, such as layer interfaces, anisotropic conduction paths, or internal air gaps in
the printed conductive shell. Such refinements were not pursued in this work, as they fall
outside the intended scope and would significantly increase computational and modeling
complexity. The objective here is not to reproduce all physical non-idealities, but to
establish a sufficiently accurate and tractable reference model for system-level analysis.

Building on the validated formulation, the same numerical framework was applied to the full
phantom head geometry. In the absence of a fully fabricated physical phantom, the simulated
head was probed in a manner analogous to the experiments envisioned for the physical system.
The resulting potential distributions and reconstruction outcomes were well behaved, with
consistently low reconstruction errors and favorable stability properties.

Within the scope of this work, the numerical phantom therefore serves as a reference system
for subsequent analyses, rather than as a validated substitute for an existing physical
implementation. The phantom--head transfer operator derived from the numerical model is
adopted as the internal ground truth against which reconstruction performance and system
behavior are assessed.

If the electrical properties of a future fabricated phantom closely match those assumed in
the model, the resulting operator is expected to provide a close approximation of the
physical system. If deviations arise, the methodology employed to construct the operator—
encompassing geometry definition, material assignment, numerical simulation, and response
extraction—can be applied directly to the physical phantom to regenerate an updated mapping.
In this sense, the robustness of the proposed approach lies not in fixed numerical values,
but in the generality and repeatability of the modeling and reconstruction pipeline.

The original system design targeted a physically realizable phantom head based on materials
and fabrication techniques that are compatible with typical academic and laboratory
resources. While the numerical and experimental results presented in this work demonstrate
the feasibility of this approach, the fabrication of the complete physical phantom head was
not carried out within the scope of the project due to practical and funding constraints.
Importantly, these constraints do not undermine the validity of the proposed architecture,
which was developed with physical realization and reproducibility as guiding objectives.

Based on the present results, the physical phantom can be fabricated following the specified
material properties and geometric design without requiring further modifications to the
modeling or control framework, once the necessary fabrication resources are available.

\textbf{TBD: Discussion of the desktop platform}

\textbf{TBD: System-level discussion of component integration and interactions}


\biblio % Needed for referencing to working when compiling individual subfiles - Do not remove
\end{document}
